Nominal Labelled Transition Systems (LTSs) can capture computational scenarios which
require unbounded sets of data values~\cite{automatanominalsets}.
In the most elementary case, data values are just \emph{names}, that is, atomic identifiers that can only be compared for equality. Nominal LTSs can be used,
for example, for modelling mobile processes~\cite{HDA1}
and computation with variables ~\cite{KAMINSKI1994329,GrumbergKS10}, and for verifying XML query languages~\cite{SCHWENTICK2007289}.
%
A particularly popular class of nominal LTSs are the ones generated by computations of ``regular'' automata over infinite alphabets, such as register automata~\cite{KAMINSKI1994329} and the (largely) equivalent nominal automata~\cite{automatanominalsets}. A register automaton is 
equipped with a finite set of states, and hence the nod to regularity, but also a finite set of registers where it can store data values and compare them with, and possibly update them with, data values from the input. This design allows register automata to capture languages over infinite alphabets even though their specification is finite.

% More recently,
% they have been applied to program and system modelling, leading to program semantics and verification tools [43, 47, 24] and, respectively, to learning algorithms for model learning [11, 1, 41]. As
% the machines that operate over infinite alphabets are so versatile, they remain an area eager to be studied.


Fresh-Register Automata (FRAs)~\cite{fra} expand on this design by having the option of allowing the automaton to accept a name just if it is globally fresh, that is, it has not appeared so far in the input.
FRAs
 can be used to capture computational models that use names and
name generation, e.g.\ (finitary) $\pi$-calculus processes~\cite{fra,DBLP:journals/jsa/BandukaraT23} or ML-like programs with mutable references~\cite{nominalgamesemantics}.
%Algorithmic nominal game semantics, Murawski + Tzevelekos?
 They have been applied to program verification and system modelling, leading to verification tools ~\cite{imj,nominalgamesemantics,runtimeverira} and model learning techniques~\cite{RCDL,LearningRAFreshval}.
Verification of FRAs has mainly focussed on bisimulation equivalence~\cite{fra,BFRA,deqtheory}. In this work we are interested in furthering formal
verification for nominal computation, with an interest on FRAs and more generally history-dependent computation, in the direction of modal logics.

%We extend nominal LTSs to capture histories fresh-register automata by designing a recursive Hennessy-Milner logic (HML) for them.

% Fresh-Register Automata are an extension of the Register Automata model by Kaminski and Francez [7]. Each automaton is equipped with a finite set of registers where it can store data values. Register automata can verify whether or not a given input data value (referred to as names) is stored in one of its registers, and store it in a register overwriting its current value. This design allows register automata to capture languages over infinite alphabets. Fresh-register automata expand on this by having the option of allowing the automaton to accept a name just if it is globally-fresh, that is, it has not appeared so far in the input. Fresh-register automata can be used to capture computational models that use names and
% name-generation, for example (finitary) $\pi$-calculus processes [16, 1]. To date, there are results for bisimulation
% equivalence of fresh-register automata [16, 13, 1]. In this work we are interested in furthering formal
% verification for fresh-register automata by designing a Hennessy-Milner logic (HML) for them.

Hennessy-Milner Logic (HML) is a modal behavioural logic that captures branching-time properties of labelled transition systems. It was initially introduced in 1980 by Hennessy and Milner~\cite{DBLP:conf/icalp/HennessyM80} as an alternative exposition of observational equivalence~\cite{BradfieldS07}. Observational equivalence has a focus on testing whether two systems can simulate each other in a step-by-step manner,
whereas HML logic focuses on the expressiveness of a single system. Under certain assumptions, two
processes are equivalent just if they satisfy the same HML formulas~\cite{JOB}. The extension of HML with
recursion, called modal $\mu$-calculus, was introduced by Kozen~\cite{Kozen83}.
Here, we study a nominal extension thereof, where modalities
are allowed to include names and where quantifications over some/all names and some/all fresh names are catered for.
More specifically, our contributions include:
\begin{itemize}
\item We introduce the notion of \emph{fresh LTS (FLTS)}, which extends nominal LTS with history dependence.
\item We introduce a \emph{fresh HML with recursion}, called \fhml, which allows us to express recursive branching-time properties involving names and name-freshness.
\item We study the model checking problem for \fhml\ formulas on FLTSs. This leads us to developing parity games for \fhml\ and, through a series of reductions and using nominal techniques, devising a decision procedure for them.
\item We provide (exponential) upper bound calculations for \fhml\ parity games and \fhml\ model checking, in the latter case both on  general FLTSs and on FRAs.
\end{itemize}

\paragraph*{Related work}
HML for the $\pi$-calculus was first examined in~\cite{MilnerPW93}. For properties involving paths
of unbounded length, it is natural to examine recursive extensions thereof such as the modal $\mu$-calculus
and variants thereof~\cite{BradfieldS07}. The logic we introduced is based on the $\pi$-$\mu$-calculus of Dam~\cite{Dam03} and is related to
the recent work on nominal $\mu$-calculi of Klin and collaborators~\cite{KlinL17,KlinL19,KlinHD}.
% Our points of focus are expressiveness of the logics with respect to fresh-register automata and previous nominal calculi, but also decidability, complexity and implementations of model checking.
The $\pi$-$\mu$-calculus~\cite{Dam03} specifically targets $\pi$-calculus processes, but the syntax used for recursion is the same as ours.
The study of a history-dependent nominal $\mu$-calculus in~\cite{KlinHD} is the work closest to ours. In a certain sense, our work continues and answers open problems from \emph{loc cit.}
 in particular by operating on history-dependent LTSs (and FRAs) and accommodating a vectorial version of recursion.
We also provide a more algorithmic view on model checking via parity games and produce complexity bounds. 

In the area of nominal techniques, the modal logics for nominal LTSs by Parrow et al.~\cite{NomLogics} can subsume our logic but do not allow for decidable model checking.
%
Logics for variants of register automata have been widely studied, e.g.~\cite{DBLP:conf/csl/Segoufin06,DBLP:journals/tocl/DemriL09,DBLP:conf/mfcs/FigueiraS09}.
% This point is taken from Bartek's paper. Refs:
% - S. Demri and R. Lazic. LTL with the freeze quantifier and register automata. ACM Trans. Comput. Log., 10(3):16:1–16:30, 2009.
% - D. Figueira and L. Segoufin. Future-looking logics on data words and trees. In Procs. MFCS 2009, volume 5734 of Lecture Notes in Computer Science, pages 331–343, 2009
% - L. Segoufin. Automata and logics for words and trees over an infinite alphabet. In Procs. CSL 2006, volume 4207 of Lecture Notes in Computer Science, pages 41–57, 2006.
The main focus of those works is on decidability, e.g.\ by restricting the number of registers allowed. No such restrictions are imposed on our logic, which is undecidable for satisfiability but our focus is instead on model checking.
There have also been a series of works on logics related to the $\pi$-calculus, e.g.~\cite{DBLP:journals/mscs/NicolaL08,DBLP:journals/cl/KoutavasH12,DBLP:conf/icalp/BergerHY08}.
Those works mainly study expressiveness and characterisation of variants of behavioural equivalence for $\pi$-calculus processes. 
Other related works, such as~\cite{DBLP:journals/iandc/Dam96,BradfieldStevens99}, expand on the modal $\mu$-calculus and study model checking, and show it to be decidable for finitary $\pi$-calculus processes. 
%Similarly, model checking of finitary $\pi$-calculus processes is decidable on \fhml\ formulas by translating to fresh-register automata~\cite{fra,DBLP:journals/jsa/BandukaraT23}.

%%% Local Variables:
%%% mode: latex
%%% TeX-master: "main"
%%% End:
